\section{A Lean-Checked Bridge to Diaconescu's Encoding}
\label{sec:diaconescu-bridge}

The main separation and completion theorems above are one-sorted and binary.
Their free-completion substrate sits inside a broader partial-to-total theory.
To make that provenance explicit, the repository contains an independent
many-sorted Lean development of the operations-only fragment of Diaconescu's
encoding \cite{Diaconescu2009}.

A many-sorted signature has sorts $S$, designated-total operation symbols
$TF$, and partial operation symbols $PF$, each with arbitrary finite arity and
typed argument and result sorts.  Encoded total models extend the data carriers
with an auxiliary truth carrier and existence-equality operations.  The Horn
conditions $\Gamma$ ensure that the defined elements form the partial reduct
$\beta(M)$.

The translation $\alpha$ sends existence-equality atoms to equality with the
truth constant and guards universal quantification by definedness.  Lean proves
the term soundness and completeness lemmas and hence the satisfaction
condition
\[
  M\models\alpha(\rho)
  \quad\Longleftrightarrow\quad
  \beta(M)\models\rho.
\]
The relevant declarations are
\leanname{Resolution.Diaconescu.ManySorted.term\_evalPartial\_sound},
\leanname{Resolution.Diaconescu.ManySorted.term\_evalPartial\_complete}, and
\leanname{Resolution.Diaconescu.ManySorted.alpha\_satisfaction\_condition}.

For every partial algebra $D$, the canonical model $\gamma(D)$ uses normalized
formal applications as its total carrier.  Lean checks $\gamma(D)\models\Gamma$
and a canonical partial-algebra equivalence
$\beta(\gamma(D))\simeq D$ through
\leanname{Resolution.Diaconescu.ManySorted.gamma\_satisfiesGamma} and
\leanname{Resolution.Diaconescu.ManySorted.betaGammaPartialAlgEquiv}.
The universal lift from a partial homomorphism into a reduct is formalized as a
hom-set equivalence; identity, composition, naturality, unit/counit behavior,
and both triangle identities are checked.

At fixed signature and universe bound, the development proves the initiality
transfers needed for the operations-only quasi-existence theories considered
here.  A conditional term-model construction supplies the initial translated
encoded model, and the $\beta$ reduct yields the corresponding initial partial
model.  Semantic consequence is preserved and reflected by the translation.
Representative declarations include
\leanname{Resolution.Diaconescu.ManySorted.gammaBetaHomEquiv},
\leanname{Resolution.Diaconescu.ManySorted.gamma\_preserves\_initiality},
\leanname{Resolution.Diaconescu.ManySorted.beta\_of\_initial\_encoded\_is\_initial},
\leanname{Resolution.Diaconescu.ManySorted.semantic\_consequence\_equivalence},
and
\leanname{Resolution.Diaconescu.InitialExistence.quasiTheory\_has\_initial\_partial\_model}.

Finally, when the signature has one sort, no designated-total symbols, and only
binary partial symbols, Lean constructs canonical equivalences between the
many-sorted generated carrier and the original binary Resolution carrier, and
between the corresponding encoded models.  The principal declarations are
\leanname{Resolution.Diaconescu.BinarySpecialization.generatedEquiv} and
\leanname{Resolution.Diaconescu.BinarySpecialization.gammaEquiv}.

\begin{remark}[Role of the bridge]
The bridge verifies that the binary implementation used by the observer theory
is a specialization of an established many-sorted partial-to-total encoding.
It is supporting formal provenance.  We do not claim Diaconescu's encoding,
the adjunction, or the existence of free partial completions as new results.
Nor does the present paper extend the finite-complement separation and
compression--escape completion theory to the full many-sorted setting; that
remains an open direction.
\end{remark}
